\section{Determinant graphs for nested highest-weight supports}

\subsection{O-015.a: arbitrary downward-closed coefficient supports}
Fix a base degree $b\geq1$ and $r\geq2$. Let
$\Gamma\subseteq\mathbb Z_{\geq0}^{r-1}$ be a nonempty finite
coordinatewise downward-closed set contained in $|\gamma|\leq b$.
Consider all forms
\[
             f(L,z)=\sum_{\gamma\in\Gamma}
                         a_\gamma z^\gamma L^{b-|\gamma|}.
\]
Put $u=\lfloor(b-1)/2\rfloor$, $v=\lceil(b-1)/2\rceil$,
$\Gamma_j=\{\gamma\in\Gamma:|\gamma|\leq j\}$, and
\[
                    D_\Gamma=|\Gamma_u|+|\Gamma_v|-1.
                                                               \tag{O-015.1}
\]
There is a homogeneous linear matrix of size $D_\Gamma$
whose determinant is the polynomial
\[
                 \sum_{\gamma\in\Gamma}
                         a_\gamma z^\gamma L^{D_\Gamma-|\gamma|}.
                                                               \tag{O-015.2}
\]
For a split $A\sqcup B=\{2,\ldots,r\}$, let $\Gamma_A$
and $\Gamma_B$ be the coordinate projections of $\Gamma$,
including the zero index. There is also such a matrix of size
\[
                         S_{\Gamma,A}=|\Gamma_A|+|\Gamma_B|-1.
                                                               \tag{O-015.3}
\]
In (O-015.2) replace $D_\Gamma$ by this size. In each case
all displayed exponents are nonnegative. Thus every
$L^{n-b}f$ in the prescribed coefficient space belongs to
$\Omega_n$ when $n\geq b$ and $n$ is at least one of the
displayed sizes, provided $r\leq n^2$.

\begin{proof}
For (O-015.1), use the two degree-split trees on
$\Gamma_u$ and $\Gamma_v$. Deleting a variable from a
multi-index stays in $\Gamma$, so all root paths exist.
Every split $\gamma=\alpha+\beta$ of a retained monomial
also has $\alpha,\beta\in\Gamma$, since each is coordinatewise
bounded by $\gamma$. For $|\gamma|<b$, use the coefficient
edge $a_\gamma L$; for $|\gamma|=b$, first remove one
variable $z_i$, split the remaining degree $b-1$, and use
$a_\gamma z_i$. These are precisely the cycles of O-014.a.
When $|\gamma|<b$ the cycle has $|\gamma|+1$ vertices;
when $|\gamma|=b$ it has $b$ vertices. Since these vertices
are distinct in the constructed graph, the corresponding
exponent in (O-015.2) is nonnegative even if the total size
is less than $b$. Permutation signs cancel all tree-edge
signs and yield (O-015.2).

For (O-015.3), the projections are downward closed, and
the projections of each $\gamma\in\Gamma$ give the two
required vertices. Use the variable-split trees on these
two sets and a coefficient edge $a_\gamma L$ between them.
Each cycle has $|\gamma|+1$ vertices, so its contribution
is $a_\gamma z^\gamma L^{S_{\Gamma,A}-|\gamma|}$, again
with nonnegative exponent. Padding either matrix by an
$L$-diagonal block of the required size gives
$L^{n-b}f$. Formally the unpadded identity can be written
$L^b\det M=L^{\operatorname{size}(M)}f$; the termwise
argument proves polynomiality even when $\operatorname{size}(M)<b$.
All identities are integral in the coefficient parameters.
\end{proof}

\subsection{O-015.b: nested support counts}
For a partition $\lambda$ with $r\geq2$ rows and
$b=\lambda_2$, take
\[
 \Gamma(\lambda)=\left\{\gamma\in\mathbb Z_{\geq0}^{r-1}:
              \sum_{i=j}^r\gamma_i\leq\lambda_j
                         \quad(2\leq j\leq r)\right\}.
                                                               \tag{O-015.4}
\]
The $j=2$ inequality bounds the total degree by $b$;
all inequalities are preserved by coordinatewise decrease.
Thus O-015.a applies. The pure second-variable powers through
degree $b$ belong to this set. Consequently $D_\Gamma\geq b$
and every $S_{\Gamma,A}\geq b+1$.
F-016 proves that this support contains the whole stabilized
highest-weight space, so either size bound preserves all
ambient determinant multiplicities at the corresponding
weight. The suffix conditions strengthen coordinatewise caps.

For $r\geq3$ write $c=\lambda_3$ and $R=r-2$.
For $0\leq j\leq b$, ignoring the later suffix inequalities
gives the useful bound
\[
 |\Gamma_j(\lambda)|\leq
 \sum_{h=0}^{\min(c,j)}(j-h+1)\binom{R+h-1}{h}.
                                                               \tag{O-015.5}
\]
Here $h$ is the total degree in variables $3,\ldots,r$;
the second variable then has $j-h+1$ possible exponents.
Equality holds if $\lambda_3=\cdots=\lambda_r=c$.
At the variable split $A=\{2\}$, $B=\{3,\ldots,r\}$,
\[
                   S_{\Gamma,A}\leq
                          b+\binom{r-2+c}{c}.          \tag{O-015.6}
\]
Again equality holds for the constant later tail. In particular,
for tails $(b,2^{r-2})$ the sufficient full-multiplicity bound
from F-016 is simply
\[
                       n\geq b+\binom r2.             \tag{O-015.7}
\]
For comparison the degree-split count in this case, for $2\leq j\leq b$,
is
\[
              |\Gamma_j|=(j+1)+j(r-2)+(j-1)\binom{r-1}{2}.
\]
The variable split can be substantially smaller, so both
constructions are retained.

\subsection{O-015.c: an actual four-copy exclusion beyond inner-degree stabilization}

Set
\[
 m=256,\quad r=65536,\quad n=m^{25}=2^{200},\quad
 d=3n,\quad b=2n,
\]
and take
\[
           \lambda=(3n^2-2n-131068,2n,2^{65534}).       \tag{O-015.8}
\]
Then the actual multiplicities satisfy
\[
                         4\leq q^m_\lambda(d[n])
                                      \leq K_n(\lambda).        \tag{O-015.9}
\]
Neither exact multiplicity is claimed here.

For the upper comparison, use M-011 at the cut $k=3$.
Although $\lambda_2=2n>n$, the retained permanent-degree
forms have total degree at most two in variables $3,\ldots,r$.
The prefix variable has degree at most $m$, so the required
matrix size is
\[
               m+\binom r2=2147451136<2^{200}=n.
\]
This proves $q\leq K_n$ without any inner-degree stabilization.
For the lower bound, M-012.3 applies with the same even $b$:
$d\geq r+b-2$ because $n\geq r-2$. It gives four
independent highest-weight classes on the permanent closure,
proving $q\geq4$. More explicitly, multiply the four binary
seed classes by
\[
              W_r\,W_2^{n-4}F(e_1)^{n-r+2}.
\]
The common factor is nonzero on the reduced irreducible
permanent closure, and the degree and weight are those in
(O-015.8). The seed classes each have degree six; the
displayed factor has degree $d-6$.

This example satisfies the BIP threshold with equality, while
$n<2^{256}-1$, so the full Grenet containment criterion does
not apply. All the stabilization-based full-multiplicity
criteria F-014--F-016 require $n\geq\lambda_2$ and fail.
For every variable split in the older M-009 box criterion,
one part has at least $32768$ variables, all with caps at
least two. Therefore
\[
                 T_m(\lambda)\geq\binom{32768}{256}>2^{200}.
\]
The quadratic Chow criterion fails since $d(d-1)>n$.
The adaptive one fails since, with $b<d$,
\[
                    B_d(b)=db-\binom b2-1
                                =4n^2+n-1>n.
\]
The coarse GCT shape conditions hold: $r=m^2$ and
$|\bar\lambda|=2n+131068\leq md$. In particular the
necessary positive-multiplicity condition $q\geq2$ is
verified rather than left unevaluated.

The accompanying exact integer certificate checks every
comparison. This certifies strict refinement of the listed
explicit sufficient conditions by M-011. It exhibits an
excluded shape with genuine multiplicity at least four; it
does not exhibit an actual obstruction or establish a value
for either of its two multiplicities.

\paragraph{Geometric dictionary.}
The downward-closed support defines an affine coefficient
subscheme selected by torus weights, in the grading/comodule
dictionary of SGA~3, I, Proposition~4.7.3, pp.~35--36.
The graph morphisms to linear matrices and their determinant
identities are defined over $\mathbb Z$. The supports need
not be stable under the entire general linear group: the
highest-weight and ideal arguments are performed in the
ambient representation, using SGA~3, I, Proposition~4.7.2,
pp.~34--35, and functorial Schur decomposition as in
Macdonald, I Appendix~A, (5.3)--(5.6), pp.~155--156.
The characteristic-zero all-copy restrictions are F-016 and
M-011. The four surviving permanent sections in the example
come from the explicit integral binary substitution of M-012,
followed by multiplication in the integral coordinate algebra
of the reduced irreducible complex orbit closure.
